\chapter{负数}\label{ch:g7:negatives}

零下的温度、地下的楼层、欠款：许多
量都会小于 $0$。负数把数轴延伸到零的
左侧。本章学习如何比较它们，以及如何加
减它们——它们的乘法留待
在\cref{ch:g8:negprod}中学习。

\section{正负数}

\begin{definition}[正负数]\label{def:g7:negatives:def}
一个\emph{正负数}\index{正负数}由符号
（$+$ 或 $-$）和\emph{到零的距离}：$+3$ 与 $-3$
的距离都是 $3$，位于 $0$ 的两侧。正数通常
省略符号书写：$+3 = 3$。$+3$ 与 $-3$ 互为
\emph{相反数}\index{相反数}。
\end{definition}

\begin{omfigure}
\begin{tikzpicture}[scale=1.15]
  \draw[-Stealth] (-5.6,0) -- (5.6,0);
  \foreach \x in {-5,...,5} \draw (\x,-0.08) -- (\x,0.08);
  \foreach \x in {-5,...,5} \node[below, font=\small] at (\x,-0.1) {$\x$};
  \fill[omThm] (-3,0) circle (2pt);
  \fill[omDef] (3,0) circle (2pt);
  \draw[Stealth-Stealth, omExo] (-3,0.35) -- (0,0.35)
    node[midway, above, font=\small] {$3$};
  \draw[Stealth-Stealth, omExo] (0,0.35) -- (3,0.35)
    node[midway, above, font=\small] {$3$};
\end{tikzpicture}

{\small 数轴向左延伸：$-3$ 与 $+3$ 互为
\omterm{def:g7:negatives:def}{相反数}，到 $0$ 的距离相等。}
\end{omfigure}

\begin{method}[比较正负数]\label{met:g7:negatives:compare}
在数轴上，越靠左的数越小。实际比较时：
\begin{enumerate}
  \item 负数总小于正数：
        $-7 < 2$；
  \item 两个正数之间，按通常的大小顺序；
  \item 两个负数中，到零的距离\emph{较大}者
        则\emph{较小}：$-7 < -2$。
\end{enumerate}
\end{method}

\begin{example}\label{ex:g7:negatives:compare}
按从小到大的顺序排列：$3$、$-5$、$0$、$-1.5$、$2.8$。
先排负数，其中最小的数排在最前面：
\[
-5 < -1.5 < 0 < 2.8 < 3 .
\]
\end{example}

\section{正负数的加法}

可以把一个\omterm{def:g7:negatives:def}{正负数}看作数轴上的一次移动：$+5$ 表示
“向右走 $5$ 步”，$-3$ 表示“向左走 $3$ 步”。做加法
就是把这些移动依次连接起来。

\begin{method}[两个正负数相加]\label{met:g7:negatives:add}
\begin{enumerate}
  \item \emph{同号}：将两数的距离相加，保留相同的符号：
        $(-3) + (-4) = -7$。
  \item \emph{异号}：用较大的距离减去较小的距离，
        取较大距离所对应数的符号：
        $(-7) + (+2) = -5$ 与 $(+7) + (-2) = +5$。
  \item 互为\omterm{def:g7:negatives:def}{相反数}的两个数相加得零：$(+4) + (-4) = 0$。
\end{enumerate}
\end{method}

\begin{omfigure}
\begin{tikzpicture}[scale=1.05]
  \draw[-Stealth] (-5.4,0) -- (4.4,0);
  \foreach \x in {-5,...,4} \draw (\x,-0.07) -- (\x,0.07);
  \foreach \x in {-5,-3,-1,0,2,4}
    \node[below, font=\small] at (\x,-0.09) {$\x$};
  \draw[-Stealth, omThm, very thick] (0,0.35) -- (-3.5,0.35);
  \draw[-Stealth, omDef, very thick] (-3.5,0.75) -- (-1.5,0.75);
  \node[omThm, font=\small] at (-1.75,0.18) {$-3.5$};
  \node[omDef, font=\small] at (-2.5,1.0) {$+2$};
  \fill (-1.5,0) circle (2pt) node[below=6pt, font=\small] {$-1.5$};
\end{tikzpicture}

{\small 把加法看作移动：从 $0$ 出发，先移动 $-3.5$，再移动
$+2$，到达 $(-3.5) + 2 = -1.5$。}
\end{omfigure}

\begin{example}\label{ex:g7:negatives:add}
逐步计算，可按方便的方式分组：
\begin{align*}
(-8) + (+3) &= -5 && \text{（异号：} 8 - 3 = 5,
\text{ 取符号 } 8)\\
(-2.5) + (-1.5) &= -4 && \text{（同号：距离相加）}\\
(+6) + (-9) + (+3) &= (-3) + (+3) = 0 .
\end{align*}
\end{example}

\section{正负数的减法}

\begin{theorem}[减法法则]\label{thm:g7:negatives:sub}
减去一个\omterm{def:g7:negatives:def}{正负数}，等同于\emph{加上它的
\omterm{def:g7:negatives:def}{相反数}}：
\[
a - b = a + (-b) .
\]
\end{theorem}

\begin{proof}[为何成立]
$a - b$ 是加上 $b$ 后得 $a$ 的数。检验
$a + (-b)$ 是否满足要求：
$\bigl(a + (-b)\bigr) + b = a + \bigl((-b) + b\bigr) = a + 0 = a$.
\end{proof}

\begin{example}\label{ex:g7:negatives:sub}
每个减法都化为加法，再
使用\cref{met:g7:negatives:add}：
\begin{align*}
5 - 9 &= 5 + (-9) = -4, \\
(-3) - (+4) &= (-3) + (-4) = -7, \\
(-2) - (-6) &= (-2) + (+6) = +4 .
\end{align*}
最后一行是最著名的结论：\emph{减去一个负数，
等于做加法}。去掉一笔债务，你就更富有了！
\end{example}

\begin{example}[温差]\label{ex:g7:negatives:temperature}
某一天内，气温从 $-6$ 度升至 $+9$ 度。
升温量即为下面的差
\[
9 - (-6) = 9 + 6 = 15 \text{ 度}.
\]
在数轴上，从 $-6$ 升到 $0$ 为 $6$ 度，从 $0$ 到
$9$ 又升了 $9$ 度。
\end{example}

\section{平面上的坐标}

\begin{definition}[带符号的坐标]\label{def:g7:negatives:coords}
在原点 $O$ 处相交的两条互相\omterm{def:g6:lines:perp}{垂直}的带刻度数轴上，
平面上每一点都对应一对\omterm{def:g7:negatives:def}{正负数}坐标 $(x, y)$：
$x$ 为负表示在纵轴左侧，$y$ 为负表示在
横轴下方。
\end{definition}

\begin{omfigure}
\begin{tikzpicture}[scale=0.62]
  \draw[gray!40, thin] (-4.4,-3.4) grid (4.4,3.4);
  \draw[-Stealth] (-4.6,0) -- (4.6,0) node[below, font=\small] {$x$};
  \draw[-Stealth] (0,-3.6) -- (0,3.6) node[left, font=\small] {$y$};
  \node[below left, font=\small] at (0,0) {$O$};
  \foreach \x in {-4,-2,2,4} \node[below, font=\footnotesize] at (\x,0)
    {$\x$};
  \foreach \y in {-2,2} \node[left, font=\footnotesize] at (0,\y) {$\y$};
  \fill[omDef] (3,2) circle (2.4pt)
    node[above right, font=\small] {$A(3, 2)$};
  \fill[omThm] (-2,3) circle (2.4pt)
    node[above left, font=\small] {$B(-2, 3)$};
  \fill[omMeth] (-3,-2) circle (2.4pt)
    node[below left, font=\small] {$C(-3, -2)$};
  \fill[omProp] (2,-3) circle (2.4pt)
    node[below right, font=\small] {$D(2, -3)$};
\end{tikzpicture}

{\small 平面的每个象限中各有一个点。
先读横坐标，总是如此。}
\end{omfigure}

\section{练习}

\begin{exercise}[$\star$]\label{exo:g7:negatives:1}
求下列各数的\omterm{def:g7:negatives:def}{相反数}：$7$；\ $-3.5$；\ $0$；\ $+12$。再补全：
“两个数的和为~？时，它们互为\omterm{def:g7:negatives:def}{相反数}”。
\end{exercise}

\begin{exercise}[$\star$]\label{exo:g7:negatives:2}
抄写并填入 $<$ 或 $>$：
\[
-4 \;?\; 1, \qquad
-4 \;?\; -7, \qquad
-2.5 \;?\; -2.4, \qquad
0 \;?\; -0.1 .
\]
\end{exercise}

\begin{exercise}[$\star$]\label{exo:g7:negatives:3}
按从小到大的顺序排列：$-3$；\ $2.5$；\ $-0.5$；\ $-3.1$；
\ $0$；\ $1$。
\end{exercise}

\begin{exercise}[$\star$]\label{exo:g7:negatives:4}
计算：
\[
(-5) + (-9), \qquad
(-12) + (+7), \qquad
(+3.5) + (-1.5), \qquad
(-6) + (+6).
\]
\end{exercise}

\begin{exercise}[$\star$]\label{exo:g7:negatives:5}
化为加法，再计算：
\[
4 - 11, \qquad
(-5) - (+3), \qquad
(-7) - (-10), \qquad
0 - (-8).
\]
\end{exercise}

\begin{exercise}[$\star$]\label{exo:g7:negatives:6}
运用巧妙分组计算（先算互为\omterm{def:g7:negatives:def}{相反数}的数）：
\[
(+7) + (-12) + (+3) + (-8), \qquad
(-5.5) + (+9) + (-4.5) + (+1).
\]
\end{exercise}

\begin{exercise}[$\star$]\label{exo:g7:negatives:7}
今天早晨气温为 $-3$ 度，白天上升了 $8$ 度，
夜间又下降了 $12$ 度。写出一个算式，
求出最终气温并计算。
\end{exercise}

\begin{exercise}[$\star$]\label{exo:g7:negatives:8}
在坐标系中描出下列各点：$A(2, 3)$；$B(-3, 1)$；$C(-2, -2)$；
$D(3, -1)$；$E(0, -3)$。哪个点在坐标轴上？
\end{exercise}

\begin{exercise}[$\star\star$]\label{exo:g7:negatives:9}
计算最高气温与最低气温之差：
莫斯科一月份，最高气温 $-4$ 度，最低气温 $-13$ 度；赫尔辛基最高
气温 $2$，最低气温 $-9$。
\end{exercise}

\begin{exercise}[$\star\star$]\label{exo:g7:negatives:10}
一名潜水员位于海拔 $-12$ m 处（即水面以下十二米）。她
先上升 $7$ m，再下降 $4$ m。她位于什么海拔？还
需上升多少米才能到达水面？
\end{exercise}

\begin{exercise}[$\star\star$]\label{exo:g7:negatives:11}
在坐标系中描出 $A(-2, 1)$、$B(2, 1)$、$C(2, -2)$。求点 $D$ 的坐标，
使 $ABCD$ 为长方形，并计算其\omterm{def:g6:measure:perimeter}{周长}
（可通过数方格求各边长）。
\end{exercise}

\begin{exercise}[$\star\star\star$]\label{exo:g7:negatives:12}
在下面的幻方中，每一行、每一列及每条对角线上的数之和
都必须相同。请复制并补全该幻方：
\begin{center}
\renewcommand{\arraystretch}{1.4}
\begin{tabular}{|c|c|c|}
\hline
$?$ & $-7$ & $0$ \\
\hline
$-3$ & $-2$ & $?$ \\
\hline
$-4$ & $?$ & $?$ \\
\hline
\end{tabular}
\end{center}
（先利用几乎
填满的那条对角线求出幻和）。
\end{exercise}

\section{问题：从未存在过的零年}

\begin{problem}[{周末问题——历史时间轴上的正负数：
跨越纪元分界的时长、天文学家的技巧与
净变化原理}]
\label{pb:g7:negatives:1}
本章的数轴在现实中有一个著名的孪生体：
即历史时间轴，公元前的年份向左延伸，
公元后的年份向右延伸。但历史学家的时间轴隐藏着一个陷阱，
它曾让许多计算失效：\emph{历史上不存在零年}
——公元前 1 年之后紧接着就是公元 1 年。本题用
气温、升降机与银行账户做计算，然后仿照天文学家的做法，
修补历史上“残缺的零”，
即使用\omterm{def:g7:negatives:def}{正负数}。

\medskip
\noindent\textbf{第一部分——热身练习，以及第一个陷阱。}
\begin{enumerate}
  \item 勃朗峰峰顶海拔为 $+4\,808$ m；
        死海岸边海拔为 $-430$ m。计算两地
        的海拔差。
  \item 一个银行账户余额为 $-45$ 欧元（透支）。其户主
        存入 $120$ 欧元，支付一笔 $60$ 欧元的账单，
        再存入 $15$ 欧元。将余额写成
        一个连加算式并计算。若只存入一笔钱，
        要使原来的 $-45$ 欧元恰好回到
        $0$ 欧元，应存入多少？
  \item 尤利乌斯·恺撒生于公元前 $100$ 年，卒于公元前 $44$ 年。
        把公元前的年份记为负数（公元前 $100$ 年记为
        $-100$，公元前 $44$ 年记为 $-44$），计算 $(-44) - (-100)$，
        恺撒共活了多少年？
  \item 奥古斯都皇帝生于公元前 $63$ 年，卒于公元
        $14$ 年。同样方法得到 $14 - (-63) = 77$ 年——
        但历史学家坚持他享年 $76$ 岁。他们的反对理由是：
        这一计数 $77$ 跨过了一个从未存在过的年份。
        是哪一年？
  \item 天文学家通过重新编号修补时间轴：公元年份
        保持原数字（$+14$ 仍为 $+14$），而公元前 $1$ 年
        记为 $0$，公元前 $2$ 年记为 $-1$，一般地，公元前
        $n$ 年记为 $-(n - 1)$。将公元前 $63$ 年换算，并
        用天文学家的数重做第 4 题的计算。
        现在的结果能否符合历史学家的看法？
\end{enumerate}

\medskip
\noindent\textbf{第二部分——跨越纪元的计算。}
从现在起，每次计算都使用天文学家的重新编号，
最后再换算回公元前、公元纪年（标号为
$0$ 或负数 $-m$ 的年份，即历史学家所说的公元前 $m + 1$ 年）。
\begin{enumerate}[resume]
  \item 换算为天文学家的数：公元前 $1$ 年；公元前 $10$ 年；公元
        $476$ 年；以及公元前 $753$ 年，即传说中罗马建城之年。
  \item 从罗马建城
        （公元前 $753$ 年）到西罗马帝国灭亡（公元 $476$ 年），共经过多少年？
  \item 一颗彗星每 $76$ 年回归一次，曾于
        公元前 $12$ 年出现。请给出它此后三次出现的历史学家纪年，
        以及公元前 $12$ 年之前那次出现的纪年。
  \item 从公元前 $5$ 年到公元 $5$ 年，共经过多少年？（
        一个诱人的答案是 $10$，但时间轴另有说法。）
  \item 一个生于公元前 $3$ 年的孩子，在哪一年满十岁
        （公元前或公元）？
\end{enumerate}

\medskip
\noindent\textbf{第三部分——净变化原理。}
\begin{enumerate}[resume]
  \item 数轴上两数之间的\emph{距离}定义为
        较大数减去较\omterm{def:g6:decimals:places}{小数}——结果总为正。
        计算 $-7$ 与 $-2$ 之间的距离，再计算
        $-3.5$ 与 $4.5$ 之间的距离。
  \item 一部电梯从 $0$ 层出发，依次运行 $+3$、
        $-5$、$+2$、$-1$、$+4$。通过巧妙分组，求出最终所在楼层
        （见\cref{met:g7:negatives:add}，
        尽可能先计算\omterm{def:g7:negatives:def}{相反数}）。若改变运行的顺序，
        最终到达的楼层会改变吗？为什么？
  \item 一周内，每天的气温变化依次为 $+2$、$-4$、$+1$、
        $0$、$-3$、$+5$、$-2$ 度。求出
        \emph{净变化}。若这一周开始时气温为 $-2$ 度，
        这一周结束时气温是多少？请叙述原理：终值 $=$
        初值 $+$（所有变化量之和）。
  \item 一位徒步者走完一段路后，恰好回到出发时的
        海拔。她所有海拔变化量之和应为多少？
        为什么？她记录的六次变化中，有五次分别为
        $+4$、$-7$、$+2$、$+3$、$-5$（以百米为单位）；
        求第六次的变化量。
  \item 压轴题，回到罗马。一位历史学家这样计算：“共和国
        从公元前 $509$ 年延续到公元前 $27$ 年，随后帝国
        延续到公元 $476$ 年；总计自公元前 $509$ 年至公元 $476$ 年，
        即 $509 + 476 = 985$ 年的罗马国家史。”用天文学家的数
        重做这一计算。正确的总时长是多少——
        又是哪一年再次导致了
        错误？
\end{enumerate}
\end{problem}
